\section{Tecnologie utilizzate}
MaaS \`e stato progettato utilizzando diverse tecnologie, alcune delle quali espressamente richieste dal proponente. Di seguito vengono elencate e descritte. Per ognuna di esse, inoltre, verranno presentanti vantaggi e svantaggi, al fine di fornire una pi\`u completa motivazione della scelta.
\begin{itemize}
\item \textbf{\glossaryItem{Node.js}}, piattaforma per il \glossaryItem{back end};
\item \textbf{ExpressJS}, \glossaryItem{framework} \glossaryItem{Node.js} flessibile e leggero che fornisce una serie di funzioni avanzate per le applicazioni web;
\item \textbf{Reactjs}, \glossaryItem{framework} \glossaryItem{JavaScript} per la realizzazione del \glossaryItem{front end};
\item \textbf{Flux}, architettura presentata da Facebook per la costituzione di frontend in \glossaryItem{JavaScript} con lo scopo di creare un flusso unidirezionale di dati;
\item \textbf{MongoDB}, database di tipo \glossaryItem{NoSQL} (\textbf{NO}t \textbf{O}nly \textbf{SQL}) per la parte di recupero e salvataggio dei dati;
\item \textbf{Mongoose}, libreria per interfacciarsi con il driver di \glossaryItem{MongoDB};
\item \textbf{HTML5} (\textbf{H}yper\textbf{T}ext \textbf{M}arkup \textbf{L}anguage), \glossaryItem{linguaggio di markup} per la creazione di pagine web;
\item \textbf{CSS3} (\textbf{C}ascading \textbf{S}tyle \textbf{S}heets), linguaggio di formattazione dei documenti \glossaryItem{HTML};
\item \textbf{Materialize}, \glossaryItem{framework} \glossaryItem{front end} per la creazione di interfacce basate sul material design;
\item \textbf{TypeScript}, super-set di \glossaryItem{JavaScript} sviluppato da Microsoft e basato su ECMAScript 6 (\textbf{E}uropean \textbf{C}omputer \textbf{M}anufacturers \textbf{A}ssociation).
\end{itemize}

\subsection{Node.js}
L'utilizzo di \glossaryItem{Node.js} \`e stato richiesto dal proponente. Si tratta di un sistema run-time \glossaryItem{cross-platform} che utilizza il motore \glossaryItem{JavaScript} V8 utilizzato anche in \glossaryItem{Google Chrome} e permette di realizzare facilmente applicazioni di rete \glossaryItem{scalabili} e veloci. Grazie al suo modello \glossaryItem{event-driven}, con chiamate di \glossaryItem{I/O} non bloccanti risulta essere leggero ed efficiente. Per l'invio di email verrà usato \textbf{Nodemailer}, un \glossaryItem{modulo} che permette di inviare facilmente e-mail con \glossaryItem{Node.js}, tramite protocollo \glossaryItem{SMTP} (\textbf{S}imple \textbf{M}ail \textbf{T}ransfer \textbf{P}rotocol). Nodemailer supporta anche la codifica \glossaryItem{unicode}.
\subsubsection{Vantaggi}
\begin{itemize}
\item \textbf{Approccio asincrono}: permette di accedere alle risorse del sistema operativo in modalit\`a \glossaryItem{event-driven} e non tramite il classico \glossaryItem{processo} orientato a \glossaryItem{processi} e \glossaryItem{thread} concorrenti utilizzato dai classici server web. Ci\`o garantisce una maggiore \glossaryItem{efficienza} dal punto di vista delle prestazioni, dato che durante le attese possono essere eseguire altre operazioni in modo asincrono.
\item \textbf{Architettura modulare}: \`e molto facile organizzare il lavoro in librerie, importare e combinare i \glossaryItem{moduli}.
\end{itemize}
\subsubsection{Svantaggi}
\begin{itemize}
\item \textbf{Orientato all'I/O}: non \`e pensato per applicazioni che sfruttano in modo intensivo la CPU (\textbf{C}entral \textbf{P}rocessing \textbf{U}nit). Questo, tuttavia, non \`e il caso di \glossaryItem{MaaS}.
\end{itemize}

\subsection{ExpressJS}
È un \glossaryItem{framework} minimale per creare applicazioni web con \glossaryItem{Node.js}. Fornisce una struttura base per un server web a cui successivamente possono essere aggiunti dei \glossaryItem{moduli} di terze parti per arricchirne le funzionalità. ExpressJS non oscura le funzionalità di base fornite da \glossaryItem{Node.js}: questo permette di utilizzare \glossaryItem{moduli} specifici a seconda dei bisogni dell'appliccazione.
\subsubsection{Vantaggi}
\begin{itemize}
\item \textbf{Minimale}: si basa su \glossaryItem{Node.js} e permette di estenderlo a seconda dei bisogni dell'applicazione.
\item \textbf{Documentazione}: esaustiva e completa.
\item \textbf{Apprendimento}: facile da imparare.
\end{itemize}
\subsubsection{Svantaggi}
\begin{itemize}
\item \textbf{Integrazione}: richiede di integrare \glossaryItem{moduli} diversi per comporre l'applicazione finale. Altri \glossaryItem{framework} permettono di definire \glossaryItem{API} (\textbf{A}pplication \textbf{P}rogramming \textbf{I}nterface) \glossaryItem{REST} (\textbf{RE}presentational \textbf{S}tate \textbf{T}ransfer) in modo semplice, ma vincolano maggiormente nelle scelte progettuali.
\end{itemize}

\subsection{Reactjs}
Uno dei requisiti opzionali espressi nel capitolato consiste nell'usare React/Flux anziché Angular per implementare un frontend per \glossaryItem{MaaS} basato su componenti riusabili.\\
React è una libreria \glossaryItem{JavaScript} usata da Facebook e Instangram per creare interfacce utente ed è considerata la View del pattern MVC.\\
Il problema che React vuole risolvere è quello di costruire applicazioni facilmente \glossaryItem{scalabili} e con dati che continuano a variare velocemente nel tempo.\\ 
L'architettura logica che usa React viene riassunta dalla seguente figura:
\begin{figure}[H]
  \centering
  \includegraphics[width=.8\textwidth]{res/sections/imgs/react.png}
  \caption{Architettura logica di React}
  \label{fig:salvataggio}
\end{figure}
\textbf{Spiegazione}: lo stato dell'applicazione viene aggiornato da eventi a partire dal Server o dal browser. Questi eventi vengono indirizzati ai componenti React, scritti in \glossaryItem{JavaScript}, i quali hanno il compito di descrivere la business logic dell'applicazione. React genera una rappresentazione virtuale del DOM che utilizza per stabilire quali operazioni eseguire con precisioni sul DOM vero e proprio; infine renderizza sul browser. L'uso di una rappresentazione virtuale del DOM rende React molto veloce poiché ad ogni cambiamento aggiorna solo le componenti interessate.\\
Disegnare la View di un'applicazione con React consiste quindi nel costruire componenti e a pensare a come incapsularli tra loro. Questo modo di lavorare produrrà codice che soddisferà il principio ``Separation of Concerns'', ne faciliterà il \glossaryItem{riuso}, il testing e, in generale, sarà più facile da manutenere.
Una componente React dovrebbe idealmente occuparsi di un solo compito, seguendo il principio ``Single Responsability''.\\
\subsubsection{Vantaggi}
\begin{itemize}
\item \textbf{Semplice}: basta esprimere come l'applicazione dovrebbe apparire in ciascun momento e React automaticamente gestir\`a tutti gli aggiornamenti dell'interfaccia utente.
\item \textbf{Dichiarativo}: quando i dati cambiano, \glossaryItem{Reactjs} effettua un refresh concettuale della pagina e aggiorna solo i dati che sono cambiati.
\item \textbf{Facilit\`a di debugging}: attraverso una specifica estensione di \glossaryItem{Google Chrome} (React Developer \glossaryItem{Tools}).
\end{itemize}
\subsubsection{Svantaggi}
\begin{itemize}
\item \textbf{Verboso}: richiede di scrivere pi\`u codice rispetto alla semplice coppia \glossaryItem{HTML} e \glossaryItem{JavaScript}.
\item \textbf{Non \`e un \glossaryItem{framework} completo}: non esiste un modello di gestione delle librerie all'interno di \glossaryItem{Reactjs}, cosa di cui altri \glossaryItem{framework} simili, come Emberjs o \glossaryItem{Angularjs}, sono provvisti.
\end{itemize}
\subsection{Flux}
Flux è l'architettura che Facebook usa per creare applicazioni web client-side. È un complemento a React e usa un flusso di dati unidirezionale. Le applicazioni Flux hanno quattro componenti principali: 
\begin{itemize}
\item il dispathcer 
\item le \glossaryItem{action} che fungono da messaggi unidirezionali tra le varie componenti
\item gli store che sono le componenti di memorizzazione
\item le viste (Reactjs).
\end{itemize}
\subsubsection{Vantaggi}
\begin{itemize}
\item \textbf{Integrazione}: si integra perfettamente con \glossaryItem{Reactjs}.
\item \textbf{Flusso di dati unidirezionale}: i dati provengono attraverso un'unica interfaccia ben definita.
\end{itemize}
\subsubsection{Svantaggi}
\begin{itemize}
\item \textbf{Apprendimento}: può essere difficile da imparare bene, soprattutto considerandolo in coppia con \glossaryItem{Reactjs}.
\end{itemize}
\subsection{MongoDB}
L'utilizzo di \glossaryItem{MongoDB} \`e stato richiesto dal proponente. Si tratta di un database \glossaryItem{NoSQL} open \glossaryItem{source} \glossaryItem{scalabile} e altamente performante di tipo \glossaryItem{document}-oriented e \glossaryItem{schemaless}, nel quale i dati sono archiviati sotto forma di documenti in stile \glossaryItem{JSON} (\textbf{J}ava\textbf{S}cript \textbf{O}bject \textbf{N}otation), con schemi dinamici e una struttura semplice e potente.
\subsubsection{Vantaggi}
\begin{itemize}
\item \textbf{Alte performance}: non ci sono \glossaryItem{join} che rallentano operazioni di lettura o scrittura. 
\item \textbf{Affidabilit\`a}: \`e presente un meccanismo di replicazione su server.
\item \textbf{Schemaless}: non esistono schemi per i dati. Pertanto \`e pi\`u flessibile.
\item \textbf{Potenza espressiva}: permette di esprimre query complesse in un linguaggio non \glossaryItem{SQL} (\textbf{S}tructured \textbf{Q}uery \textbf{L}anguage).
\item \textbf{Flessibilit\`a}: per i tipi di dato.
\end{itemize}
\subsubsection{Svantaggi}
\begin{itemize}
\item \textbf{Nessun supporto per le transazioni}: sono supportate alcune operazioni atomiche, ma a livello di documento.
\item \textbf{Nessun \glossaryItem{join}}: va simulato via codice attraverso query multiple.
\item \textbf{Utilizzo di memoria}: maggiore rispetto ai database \glossaryItem{SQL}; infatti memorizza i nomi delle chiavi in ogni documento. 
\item \textbf{Problemi di concorrenza}: per le operazioni di scrittura viene creato un lock sull'intero database. Questo lock blocca anche le operazioni di lettura.
\end{itemize}
\subsection{Mongoose}
Mongoose \`e una libreria per interfacciarsi a \glossaryItem{MongoDB} che permette di definire degli schemi per modellare i dati del database. Inoltre fornisce strumenti utili per la \glossaryItem{validazioni} dei dati, per la definizione di queries e per il cast dei tipi predefiniti. \\
\subsubsection{Vantaggi}
\begin{itemize}
\item \textbf{Diffusione}: \`e la libreria pi\`u diffusa per interfacciarsi con \glossaryItem{MongoDB}.
\item \textbf{Funzionalit\`a aggiuntive}: permette di definire strumenti per la \glossaryItem{validazione} dei dati e per il cast dei tipi.
\end{itemize}
\subsubsection{Svantaggi}
\begin{itemize}
\item \textbf{Schema-based}: \`e basato sulla creazione di una forte schematizzazione per i documenti, e questo limità l'estrema flessibilità di \glossaryItem{MongoDB}.
\end{itemize}
\subsection{HTML5}
\`E un \glossaryItem{linguaggio di markup} per la strutturazione delle pagine web, pubblicato come \glossaryItem{W3C} Recommendation da ottoble 2014. L'uso di HTML5 rispetto a \glossaryItem{XHTML} (e\textbf{X}tensible \textbf{H}yper\textbf{T}ext \textbf{M}arkup \textbf{L}anguage) \`e stato deciso all'unanimit\`a dal gruppo (vedi decisione \textit{RI-1403-2}) e approvato dal proponente.
\subsubsection{Vantaggi}
\begin{itemize}
\item \textbf{Raccomandazione \glossaryItem{W3C}} (\textbf{W}orld \textbf{W}ide \textbf{W}eb \textbf{C}onsortium).
\item \textbf{Creazione di pagine interattive}: soprattutto se usato insieme a \glossaryItem{CSS}.
\end{itemize}
\subsubsection{Svantaggi}
\begin{itemize}
\item \textbf{Supporto}: non tutti i browser lo supportano allo stesso modo, e non tutte le caratteristiche definite sono ancora completamente supportate.
\end{itemize}
\subsection{CSS3}
\`E un linguaggio utilizzato per definire la formattazione di documenti \glossaryItem{HTML} e \glossaryItem{XHTML}. Le regole per la composizione di un foglio di stile \glossaryItem{CSS} sono definite dal \glossaryItem{W3C} a partire dal 1996. Inoltre permette di separare i contenuti delle pagine \glossaryItem{HTML} dalla loro formattazione, assicurando una maggiore manutenibilit\`a e riutilizzo.
\subsubsection{Vantaggi}
\begin{itemize}
\item \textbf{Separazione tra contenuto e presentazione}.
\item \textbf{Raccomandazione \glossaryItem{W3C}}.
\end{itemize}
\subsection{Materialize}
Creato e progettato da Google, il material design è un linguaggio di progettazione di interfacce con l'obiettivo di unificare l'esperienza utente attraverso le diverse piattaforme.
\subsubsection{Vantaggi}
\begin{itemize}
\item \textbf{Open \glossaryItem{Source}};
\item \textbf{Completo}: permette una vasta selezione di componenti grafici.
\end{itemize}
\subsubsection{Svantaggi}
\begin{itemize}
\item \textbf{Flexbox}: non usa il modello Flexbox, uno dei nuovi layut introdotti da CSS3.
\end{itemize}
\subsection{TypeScript}
Aggiunge a \glossaryItem{JavaScript} tipi statici (opzionale) e meccanismi propri della progrmmazione orientata agli oggetti quali classi, interfacce ed ereditariet\`a. Attraverso un \glossaryItem{processo} di compilazione viene prodotto codice \glossaryItem{JavaScript} leggibile e conforme agli standard di codifica. 
\subsubsection{Vantaggi}
\begin{itemize}
\item \textbf{Fail Fast Principle}, gli errori di battitura e di tipo vengono scoperti a tempo di compilazione e non a runtime.
\item \textbf{Supporto ai \glossaryItem{moduli}}.
\item \textbf{Supporto alla \glossaryItem{programmazione} orientata agli oggetti}.
\item \textbf{Controllo di tipo:} TypeScript fornisce anche un meccanismo di controllo di tipo a livello di compilazione.
\end{itemize}
\subsubsection{Svantaggi}
\begin{itemize}
\item \textbf{Perdita di flessibilit\`a}, con l'introduzione di tipi statici.
\item \textbf{Compilazione}, per eseguire l'applicazione \`e necessaria la compilazione dei file TypeScript.
\item \textbf{Supporto}, molti \glossaryItem{framework} supportano file \glossaryItem{JavaScript}, ma ancora non file TypeScript.
\end{itemize}
